% Un template bilingue pour la production des m�moires et th�ses dans le d�partement de Math�matiques-Informatique.
% Ce template est conforme aux recommandations de l'�cole doctorale
%
% Ce fichier est con�u pour accueillir vos imports (\usepackage) et vos propres d�finitions d'environnements latex et/ou de style
% Consulter le fichier style.tex pour savoir ce qui a d�ja �t� import�
%
% @author Zekeng Ndadji Milliam Maxime

% M�tadonn�es m�moire (auteur, titre, encadreur)
% Metadonnees du memoire — page de garde, certification, titre anglais
% Auteur : DNJOMOU YONMBA Wilfried Loic

\newcommand{\thesisAuthor}{DNJOMOU YONMBA Wilfried Lo\"ic}
\newcommand{\thesisMatricule}{CM-UDS-24SCI0999}
\newcommand{\thesisPriorDegree}{Licence de technologie en conception, d\'eveloppement, r\'eseaux et Internet}
\newcommand{\thesisOption}{Informatique Fondamentale}
\newcommand{\thesisSpecialization}{Intelligence Artificielle}
\newcommand{\thesisSupervisor}{Pr BOMGNI Alain Bertrand}
\newcommand{\thesisSupervisorTitle}{Professeur Titulaire, Universit\'e de Dschang}
\newcommand{\thesisAcademicYear}{2025/2026}

% Titre principal (francais — page de garde)
\newcommand{\thesisTitleFr}{%
  Pr\'ediction par intelligence artificielle des contaminants \'emergents des eaux~:%
  \\[0.4em] \'etude de cas sur la d\'etection des PFAS et l'\'evaluation des risques%
}

% Titre court (certification, resume — une ligne)
\newcommand{\thesisTitleFrShort}{%
  Pr\'ediction par intelligence artificielle des contaminants \'emergents des eaux~: \'etude de cas sur la d\'etection des PFAS et l'\'evaluation des risques%
}

% Titre seconde langue (anglais)
\newcommand{\thesisTitleEn}{%
  AI-Powered Prediction of Emerging Water Contaminants:\\
  A Case Study in PFAS Detection and Risk Assessment%
}


\usepackage[english]{algorithm2e} %Pour �crire des algorithmes
%\SetKwIF{Si}{SinonSi}{Sinon}{si}{alors}{sinon si}{sinon}{finsi} %pour franciser le If
%\SetKw{Debut}{Fin}
\SetKwFor{For}{for}{do}{endfor}
%\SetKwFor{PourTout}{pourTout}{faire}{finPour}%pour franciser le pour (on a remplacer ''pour'' par ''pourTout'' 
%\SetKwForAll{PourTout}{pourTout}{faire}{finPour}
\SetKwRepeat{Repeat}{repeat}{until}
%\SetKwRepeat{Repeter}{repeter}{jusqu'�}%pour franciser le ''repeter''
%\restylealgo{boxed}\linesnumbered %pour encader les algorithmes et num�roter ses lignes
%\SetKwFor{Tantque}{tantque}{faire}{fintq}%pour franciser le tant que


%D�finition de nouveaux environnements de type th�or�me
\newtheorem{theorem}{Th�or�me}
\newtheorem{definition}[theorem]{D�finition}
\newtheorem{proposition}[theorem]{Proposition}
\newtheorem{lemma}[theorem]{Lemme}
\newtheorem{example}[theorem]{Exemple}
\newtheorem{remark}[theorem]{Remarque}
\newtheorem{corollary}[theorem]{Corollaire}
\newtheorem{problem}[theorem]{Probl�me}

%Environnements de preuve
\newenvironment{proof}[1][{\textbf{Proof}}]{
	\par
	\normalfont
	\topsep6\p@\@plus6\p@ \trivlist
	\item[\hskip\labelsep\itshape
	#1\@addpunct{.}]\ignorespaces
}{%
	\qed\endtrivlist
}
\newenvironment{preuve}[1][{\textbf{Preuve}}]{
	\par
	\normalfont
	\topsep6\p@\@plus6\p@ \trivlist
	\item[\hskip\labelsep\itshape
	#1\@addpunct{.}]\ignorespaces
}{%
	\qed\endtrivlist
}

\usepackage{float}
\usepackage{longtable} % tableaux secables sur plusieurs pages (liste des acronymes)
\usepackage{tikz}
\usetikzlibrary{arrows.meta,positioning,shapes.geometric,fit,shadings,calc}
\usepackage{eso-pic}
\usepackage{pdfpages}

%\usepackage{geometry}
%\geometry{hmargin={4.5cm,-1cm},vmargin={5.5cm,0.25cm}}
